
\documentclass[11pt,a4paper]{article}
\usepackage[T1]{fontenc}
\usepackage{lmodern}
\usepackage{amsmath,amssymb,amsfonts,mathtools}
\usepackage{graphicx,booktabs,longtable,array,tabularx,adjustbox}
\usepackage{geometry,microtype,fancyhdr,lastpage,etoolbox,enumitem}
\usepackage{newunicodechar,textcomp}
\usepackage{hyperref}
\hypersetup{hidelinks,pdfcreator={XeLaTeX / QM1 academic edition}}
\geometry{left=20mm,right=20mm,top=22mm,bottom=23mm}
\setcounter{secnumdepth}{2}
\setlength{\parindent}{0pt}
\setlength{\parskip}{4pt plus 1pt}
\setlength{\emergencystretch}{2em}
\setlength{\tabcolsep}{3pt}
\renewcommand{\arraystretch}{1.16}
\setlist{nosep,topsep=3pt,itemsep=2pt,leftmargin=7mm}
\AtBeginEnvironment{longtable}{\footnotesize\setlength{\tabcolsep}{2.2pt}}
\newcommand{\tightlist}{\setlength{\itemsep}{0pt}\setlength{\parskip}{0pt}}
\providecommand{\Cov}{\operatorname{Cov}}
\providecommand{\Var}{\operatorname{Var}}
\providecommand{\Corr}{\operatorname{Corr}}
\providecommand{\MAD}{\operatorname{MAD}}
\providecommand{\E}{\mathbb{E}}
\providecommand{\R}{\mathbb{R}}
\newunicodechar{—}{\textemdash}
\newunicodechar{–}{\textendash}
\newunicodechar{€}{\texteuro}
\newunicodechar{≤}{\ensuremath{\leq}}
\newunicodechar{≥}{\ensuremath{\geq}}
\newunicodechar{×}{\ensuremath{\times}}
\newunicodechar{÷}{\ensuremath{\div}}
\newunicodechar{Σ}{\ensuremath{\Sigma}}
\newunicodechar{σ}{\ensuremath{\sigma}}
\newunicodechar{μ}{\ensuremath{\mu}}
\newunicodechar{Δ}{\ensuremath{\Delta}}
\newunicodechar{∑}{\ensuremath{\sum}}
\newunicodechar{≈}{\ensuremath{\approx}}
\newunicodechar{√}{\ensuremath{\surd}}
\newunicodechar{≠}{\ensuremath{\ne}}
\newunicodechar{→}{\ensuremath{\to}}
\newunicodechar{∞}{\ensuremath{\infty}}
\newunicodechar{∈}{\ensuremath{\in}}
\newunicodechar{−}{\ensuremath{-}}
\newunicodechar{²}{\ensuremath{^{2}}}
\newunicodechar{³}{\ensuremath{^{3}}}
\newunicodechar{₀}{\ensuremath{_{0}}}
\newunicodechar{₁}{\ensuremath{_{1}}}
\newunicodechar{₂}{\ensuremath{_{2}}}
\newunicodechar{₃}{\ensuremath{_{3}}}
\newunicodechar{₄}{\ensuremath{_{4}}}
\newunicodechar{ₜ}{\ensuremath{_{t}}}
\newunicodechar{ₚ}{\ensuremath{_{p}}}
\newunicodechar{ₙ}{\ensuremath{_{n}}}
\newunicodechar{ₓ}{\ensuremath{_{x}}}
\newunicodechar{ᵢ}{\ensuremath{_{i}}}
\newunicodechar{ŷ}{\ensuremath{\hat y}}
\newunicodechar{Ŷ}{\ensuremath{\widehat Y}}
\newunicodechar{ₘ}{\ensuremath{_{m}}}
\newunicodechar{ₑ}{\ensuremath{_{e}}}
\newunicodechar{ₐ}{\ensuremath{_{a}}}
\newunicodechar{ₖ}{\ensuremath{_{k}}}
\newunicodechar{ₛ}{\ensuremath{_{s}}}
\newunicodechar{ₗ}{\ensuremath{_{l}}}
\newunicodechar{ₕ}{\ensuremath{_{h}}}
\newunicodechar{ₒ}{\ensuremath{_{o}}}
\newunicodechar{₊}{\ensuremath{_{+}}}
\newunicodechar{₋}{\ensuremath{_{-}}}
\newunicodechar{₌}{\ensuremath{_{=}}}
\newunicodechar{₍}{\ensuremath{_{(}}}
\newunicodechar{₎}{\ensuremath{_{)}}}
\newunicodechar{⁻}{\ensuremath{^{-}}}
\newunicodechar{⁺}{\ensuremath{^{+}}}
\newunicodechar{⁰}{\ensuremath{^{0}}}
\newunicodechar{¹}{\ensuremath{^{1}}}
\newunicodechar{⁴}{\ensuremath{^{4}}}
\newunicodechar{⁵}{\ensuremath{^{5}}}
\newunicodechar{⁶}{\ensuremath{^{6}}}
\newunicodechar{⁷}{\ensuremath{^{7}}}
\newunicodechar{⁸}{\ensuremath{^{8}}}
\newunicodechar{⁹}{\ensuremath{^{9}}}
\makeatletter
\renewcommand{\section}{\@startsection{section}{1}{0pt}{2.0ex plus .7ex minus .3ex}{1.0ex}{\large\bfseries}}
\renewcommand{\subsection}{\@startsection{subsection}{2}{0pt}{1.8ex plus .5ex minus .2ex}{.7ex}{\normalsize\bfseries}}
\renewcommand{\subsubsection}{\@startsection{subsubsection}{3}{0pt}{1.3ex plus .4ex minus .2ex}{.5ex}{\normalsize\bfseries}}
\makeatother
\pagestyle{fancy}
\fancyhf{}
\renewcommand{\headrulewidth}{0.35pt}
\renewcommand{\footrulewidth}{0pt}
\fancyhead[L]{\small\textsc{KEDGE Business School --- QM1}}
\fancyhead[R]{\small\textsc{2026--2027}}
\fancyfoot[L]{\small Badr ABOUFARASSE, PhD}
\fancyfoot[R]{\small\thepage\ /\ \pageref{LastPage}}
\setlength{\headheight}{14pt}
\setkeys{Gin}{keepaspectratio}
\begin{document}

\begin{center}
{\small\textsc{KEDGE Business School \quad EBP1 \quad Quantitative Methods 1}}\\[6pt]
{\Large\bfseries Worksheet 7 — Integrated cases}\\[6pt]
{\normalsize\bfseries Exercise statements}\\[6pt]
{\small Badr ABOUFARASSE, PhD \quad --- \quad 2026--2027}
\end{center}
\vspace{2mm}\hrule\vspace{6mm}

\hypertarget{assessment-and-exam-practice}{%
\section{Assessment and exam practice}\label{assessment-and-exam-practice}}

Practise independently with assessment and final-exam style questions.

\emph{Assessment / exam practice --- 70 min}

\hypertarget{case-a-continuous-assessment-preparation}{%
\subsection{Case A -- Continuous-assessment preparation}\label{case-a-continuous-assessment-preparation}}

A school surveys 200 students about the maximum price they would accept for an optional service:

\begin{longtable}[]{@{}lllllll@{}}
\toprule\noalign{}
Price & {\([5,10[\)} & {\([10,15[\)} & {\([15,20[\)} & {\([20,30[\)} & {\([30,40[\)} & {\([40,60[\)} \\
\midrule\noalign{}
\endhead
\bottomrule\noalign{}
\endlastfoot
{\(n_i\)} & 10 & 25 & 45 & 60 & 40 & 20 \\
\end{longtable}

Identify population, unit and variable type; build the complete table with widths, densities, frequencies, ICC, DCC and midpoints; draw histogram and cumulative curves; calculate modal class, median, quartiles, mean, population variance, standard deviation and CV; find the price affordable to about 65\%; then classify a satisfied/not satisfied variable and choose a graph.

\emph{Assessment / exam practice --- 110 min}

\hypertarget{case-b-new-mock-final-exam-110-minutes}{%
\subsection{Case B --- New mock final exam --- 110 minutes}\label{case-b-new-mock-final-exam-110-minutes}}

Independent mock exam: the data below differ from exercises in chapters 4--6. Show every calculation step and justify rounding.

\hypertarget{part-1-two-regressions-40-min}{%
\subsubsection{Part 1 --- Two regressions (40 min)}\label{part-1-two-regressions-40-min}}

\begin{longtable}[]{@{}lllllll@{}}
\toprule\noalign{}
t & 1 & 2 & 3 & 4 & 5 & 6 \\
\midrule\noalign{}
\endhead
\bottomrule\noalign{}
\endlastfoot
x & 5 & 6 & 8 & 9 & 11 & 13 \\
y & 12 & 15 & 18 & 20 & 24 & 29 \\
\end{longtable}

1. Build t, x, y, t², x², y², tx, xy table; calculate both correlations and the regressions x on t and y on x. 2. Forecast x and y for t=8; qualify the forecast.

\hypertarget{part-2-quarterly-series-40-min}{%
\subsubsection{Part 2 --- Quarterly series (40 min)}\label{part-2-quarterly-series-40-min}}

\begin{longtable}[]{@{}lllll@{}}
\toprule\noalign{}
Année / Year & Q1 & Q2 & Q3 & Q4 \\
\midrule\noalign{}
\endhead
\bottomrule\noalign{}
\endlastfoot
A & 72 & 88 & 80 & 110 \\
B & 78 & 96 & 88 & 121 \\
\end{longtable}

1. Plot the series; calculate MA4 and CMA4. 2. With the supplied trend Tₜ=2.8t+78, compute all eight ratios Yₜ/Tₜ. 3. Average by quarter, normalise (sum 4) and forecast t=9 to 12.

\hypertarget{part-3-index-numbers-30-min}{%
\subsubsection{Part 3 --- Index numbers (30 min)}\label{part-3-index-numbers-30-min}}

\begin{longtable}[]{@{}lllll@{}}
\toprule\noalign{}
Produit / Product & p₀ & q₀ & p₁ & q₁ \\
\midrule\noalign{}
\endhead
\bottomrule\noalign{}
\endlastfoot
A & 10 & 40 & 12 & 45 \\
B & 8 & 60 & 9 & 55 \\
C & 20 & 20 & 22 & 25 \\
\end{longtable}

Compute all four aggregates; Laspeyres, Paasche and Fisher price and quantity indices; the value index; and verify Fisher\textquotesingle s factor reversal.

\end{document}
